%\documentclass[portrait,a0,draft]{a0poster}
\documentclass[portrait,a0]{a0poster}
%\usepackage{german,epsf,psfig,pstricks,multicol,pst-grad,color}
\usepackage{german,pstricks,multicol,pst-grad,color}
\usepackage{graphicx}
\usepackage[latin1]{inputenc}
\usepackage{times}
\usepackage{helvetic}

\begin{document}

% Fliessenden Hintergrund von RGB-Farbe 1. .98 .98 nach 1. .85 .85 
% und wieder nach  1. .98 .98 (1. .85 .85 wird nach 0.1=10% des Hinter-
% grunds angenommen)
% Achtung Werte unter .8 verbrauchen zu viel Tinte!!!
%\background{1. .98 .98}{1. .85 .85}{0.1}
% Hier haben wir ein paar Experimente gemacht und uns schliesslich auf
\background{.98 .98 1.}{0.698 .819 0.941}{0.05}
% geeinigt. Andere Versuche:
%\background{.98 .98 1.}{.0 .39 0.65}{0.1}
%\background{.98 .98 1.}{.0 .57 0.95}{0.1}
%\background{.98 .98 1.}{0.741 .870 0.941}{0.05}
%\background{.85 .85 1.}{.98 .98 1.}{-0.1}

% Groesse der einzelnen Spalten als Anteil der Gesamt-Textbreite
\renewcommand{\columnfrac}{.3}


% Posterueberschrift, optional mit RGB-Hintergrundfarbe (hier .99 .99 1.)
\begin{header}[1. .98 .98]
\begin{tabular}{|p{9.5in}|p{19.5in}|}
\hline
\begin{minipage}[t]{9.5in}
 \vspace*{0mm}
 \begin{center}
 \includegraphics[width=7in]{rwthaachenlogo.eps}
 \end{center}

 \vspace*{-0.8ex}
 \rule{9.5in}{0.05ex}

 \vspace*{0.6ex}
 \begin{tabular}{cp{6.5in}}
 \sf\veryHuge \raisebox{-0.35ex}{LDfM} &
 \renewcommand{\baselinestretch}{0.8}\sf\Large
  \hspace*{0.5in} Lehrstuhl D f�r Mathematik
 \end{tabular}
 \renewcommand{\baselinestretch}{1}
 \vspace*{0.3ex}
\end{minipage}
&
\begin{minipage}[t]{19.5in}
 \vspace*{0mm}
\begin{center}
\vspace*{2ex}
\Huge OKUSON - ein System zur Organisation des �bungsbetriebs zu gro�en
Vorlesungen

\vspace*{0.8ex}
\Large Frank L�beck und Max Neunh�ffer
\end{center}
\end{minipage} \\
\hline
\end{tabular}

\end{header}

\begin{poster}
%
% Erste Spalte:
%
\begin{pcolumn}
\begin{pbox}
\section{Motivation}

In Deutschland gibt es, im Gegensatz zu den meisten anderen L�ndern, an
Universit�ten recht viele Vorlesungen mit sehr vielen H�rern, zum Teil
etliche hundert bis �ber 1000. Meist wird dazu ein �bungsbetrieb in
kleinen oder zumindest kleineren Gruppen angeboten, bei dem regelm��ig
Hausaufgaben gestellt, eingesammelt, korrigiert und besprochen werden.

Die Organisation und Durchf�hrung solch eines �bungsbetriebs ist aufwendig:
zum Einen werden viele Tutoren ben�tigt, um die Hausaufgaben zu korrigieren
und in den �bungsstunden zu besprechen, zum Anderen stellt schon die
Organisation selbst eine gewissen Herausforderung dar, da die Aufgaben
verteilt und die Ergebnisse erfasst und verarbeitet werden m�ssen.

Das System {\sf OKUSON} ist
\end{pbox}
\begin{pbox}
\section{Was kann OKUSON?}
\end{pbox}
\begin{pbox}[1. .98 .98]
\section{Aufgabentypen}
\vspace{30cm}
\end{pbox}
\end{pcolumn}
%
% Zweite Spalte:
%
\begin{pcolumn}
\begin{pboxns}
\section{Erfassen von Aufgaben}
%\vspace{30cm}
\end{pboxns}
\begin{pbox}
\section{Implementation}
%\vspace{30cm}
\end{pbox}
\begin{pbox}
\section{Eingebauter Webserver}
%\vspace{30cm}
\end{pbox}
\end{pcolumn}
%
% Dritte Spalte:
%
\begin{pcolumn}
\begin{pboxns}
\section{Anwendungsbeispiele}
%\vspace{30cm}
\end{pboxns}
\begin{pbox}
\section{Screenshots}
%\vspace{30cm}
\end{pbox}
\begin{pbox}
\section{Lizenz}

{\sf OKUSON} ist freie Software. Es steht unter dem Copyright von Frank 
L�beck und Max Neunh�ffer und kann unter den Bedingungen der GNU Public
license, Version 2 oder h�her, lizensiert werden.

\smallskip
Die offizielle Webseite von {\sf OKUSON} ist

\bigskip
\hspace*{1cm} \verb!http://www.math.rwth-aachen.de/~OKUSON!
\end{pbox}
\end{pcolumn}
%
\end{poster}

\end{document}

